% Generated draft: requires mathematical and source review.
\documentclass[11pt]{article}
\usepackage[margin=25mm]{geometry}
\usepackage{fontspec}
\setmainfont{Latin Modern Roman}
\usepackage{amsmath,amssymb,mathtools,unicode-math}
\setmathfont{Latin Modern Math}
\usepackage{enumitem,needspace,xurl,hyperref}
\setlength{\parindent}{0pt}
\setlength{\parskip}{5pt}
\setlength{\emergencystretch}{3em}
\newcommand{\N}{\mathbb{N}}
\newcommand{\Z}{\mathbb{Z}}
\newcommand{\Q}{\mathbb{Q}}
\newcommand{\R}{\mathbb{R}}
\newcommand{\C}{\mathbb{C}}
\newcommand{\F}{\mathbb{F}}
\newcommand{\A}{\mathbb{A}}
\newcommand{\PP}{\mathbb{P}}
\newcommand{\E}{\mathbb{E}}
\newcommand{\eps}{\varepsilon}
\newcommand{\dd}{\,\mathrm{d}}
\begin{document}
\section*{UnsolvedMath ID: 20000802}
\textbf{Source problem code:} catalogue entry
\subsection{Definitions and mathematical statement}
Let \(k\) be an algebraically closed field of characteristic zero, the setting of the supplied partial result, and let \(D\ge1\), \(n\ge2\). Write
\[
 H_D^n=\operatorname{Hilb}^{D}(\A_k^n),\qquad
 \mathcal R_D^n=\overline{\{\text{unions of }D\text{ distinct points}\}}_{\mathrm{red}}
\]
for the Hilbert scheme of length-\(D\) subschemes and its reduced smoothable component. For a point \([Z]\) with ideal \(I_Z\subset k[x_1,\ldots,x_n]\), distinguish
\[
 T_{[Z]}\mathcal R_D^n\ \subseteq\ T_{[Z]}H_D^n
 =\operatorname{Hom}_{k[x_1,\ldots,x_n]}(I_Z,k[x_1,\ldots,x_n]/I_Z).
\]
The source asks to describe the singularities of the smoothable component, how large its Zariski tangent dimension can be, and whether the maximum occurs at intersections of irreducible components. One precise numerical target for the component itself is
\[
 M_{\mathrm{sm}}(D,n)=\max_{[Z]\in\mathcal R_D^n(k)}\dim_kT_{[Z]}\mathcal R_D^n.
\]
This differs from maximizing the ambient Hilbert-scheme tangent dimension along \(\mathcal R_D^n\), or over all of \(H_D^n\). The source's final question, ``Does it occur in the smoothable component?'', is verified verbatim in the supplied verification and is redundant if ``it'' refers to the maximum of tangent dimensions to that component. It might instead ask whether an ambient maximum lies on the smoothable component, or whether a maximum occurs away from all other components. No one of these interpretations is selected as the corrected source question.

The supplied aid reports the following witness family, not a determination of the maximum. Write
\[
 D=r(n+1)+s,\qquad r=\lfloor D/(n+1)\rfloor,\qquad0\le s\le n.
\]
Choose pairwise distinct points and take the disjoint union \(Z\) of \(r\) square-fat points \(\operatorname{Spec}(k[x_1,\ldots,x_n]/\mathfrak m_a^2)\), each of length \(n+1\), and \(s\) reduced points. The aid reports that \([Z]\in\mathcal R_D^n\) and that the above tangent-space inclusion is an equality, with
\[
 \dim_kT_{[Z]}\mathcal R_D^n
 =\dim_kT_{[Z]}H_D^n
 =nD+r\frac{n(n+1)(n-2)}2.
\]
Consequently its reported lower bound is
\[
 M_{\mathrm{sm}}(D,n)\ge
 nD+\left\lfloor\frac D{n+1}\right\rfloor\frac{n(n+1)(n-2)}2.
\]
For \(n\ge7\) and \(r\ge1\), it further reports, using Lee's conical symmetric chart, that these witnesses lie on the intersection of distinct components of \(H_D^n\). This does not assert that any witness realizes the unknown maximum, nor that all maxima occur at component intersections. The ground-field and dimension restrictions of this partial result are not explicitly imposed in the original extracted problem.
\subsection{Short English statement}
The problem asks how singular the smoothable component of the Hilbert scheme of points can be. The supplied aid gives explicit large tangent spaces and some witnesses on component intersections, while leaving the exact maximum and the source's ambiguous final question unresolved.
\subsection{Sources}
\begin{itemize}
\item Local UnsolvedMath registry, permanent ID \(20000802\), source code \texttt{catalogue entry}; \texttt{aim-arithmetic-geometry-notes.json}, \texttt{notes[47]}, tag \texttt{problem}.
\item AIM workshop ``Components of Hilbert schemes,'' Problem 1: \url{https://aimath.org/WWN/hilbertschemes/hilbertschemes.pdf}. Statement status \texttt{exact}; the supplied verification says the final sentence was checked verbatim and is not an extraction error. No PDF was accessed for this draft.
\item Lee's conical symmetric chart is named in the research summary without a title, theorem number, or full authoritative citation. No original literature references, remarks, proposer, or year are supplied.
\item Research attempt 1: \texttt{partial\_result}, problem status at run \texttt{open}, rigor label \texttt{proved}, novelty confidence low; registry status \texttt{partially\_solved}, classification \texttt{PARTIAL-PROGRESS}. Created and updated 14 August 2026. The machine-generated aid is not peer reviewed and must be independently verified before citation.
\item Ulam AI / UnsolvedMath Contributors; CC BY 4.0.
\end{itemize}
\end{document}
